% !TeX root = ../sections/02_exercises.tex
\subsection{2-B-1.}
\begin{exercise}
	上に有界な集合 \(A,B\subset\mathbb{R}\) に対して，
	\[
	C:=\{a+b\mid a\in A,\ b\in B\}
	\]
	とする。
	このとき，
	\[
	\sup C=\sup A+\sup B
	\]
	が成立することを証明せよ。
\end{exercise}

\begin{background}
	この問題では，上限の定義を使う。
	実数 \(s\) が集合 \(A\) の上限であるとは，
	次の二つが成り立つことである。
	
	\[
	a\leq s
	\qquad (a\in A)
	\]
	であり，さらに任意の \(\varepsilon>0\) に対して，
	ある \(a_\varepsilon\in A\) が存在して
	\[
	s-\varepsilon<a_\varepsilon
	\]
	となる。
	
	ここでは
	\[
	\alpha=\sup A,
	\qquad
	\beta=\sup B
	\]
	とおくと，
	任意の \(a\in A,\ b\in B\) に対して
	\[
	a\leq\alpha,
	\qquad
	b\leq\beta
	\]
	である。
	したがって
	\[
	a+b\leq\alpha+\beta
	\]
	となり，\(\alpha+\beta\) は \(C\) の上界になる。
	
	あとは，\(\alpha+\beta\) がただの上界ではなく，
	最小の上界であることを示せばよい。
\end{background}

\begin{strategy}
	\[
	\alpha=\sup A,
	\qquad
	\beta=\sup B
	\]
	とおく。
	示すべきことは
	\[
	\sup C=\alpha+\beta
	\]
	である。
	
	まず，\(\alpha+\beta\) が \(C\) の上界であることを示す。
	任意に \(c\in C\) を取ると，
	ある \(a\in A,\ b\in B\) が存在して
	\[
	c=a+b
	\]
	と書ける。
	\(\alpha,\beta\) はそれぞれ \(A,B\) の上界なので，
	\[
	a\leq\alpha,
	\qquad
	b\leq\beta
	\]
	である。
	したがって
	\[
	c=a+b\leq\alpha+\beta
	\]
	となる。
	よって \(\alpha+\beta\) は \(C\) の上界である。
	
	次に，\(\alpha+\beta\) より小さい数は
	\(C\) の上界ではないことを示す。
	任意に \(\varepsilon>0\) を取る。
	\(\alpha=\sup A\) より，ある \(a_\varepsilon\in A\) が存在して
	\[
	\alpha-\frac{\varepsilon}{2}<a_\varepsilon
	\]
	となる。
	また，\(\beta=\sup B\) より，ある \(b_\varepsilon\in B\) が存在して
	\[
	\beta-\frac{\varepsilon}{2}<b_\varepsilon
	\]
	となる。
	
	このとき \(a_\varepsilon+b_\varepsilon\in C\) であり，
	\[
	\alpha+\beta-\varepsilon
	<
	a_\varepsilon+b_\varepsilon
	\]
	が成り立つ。
	したがって，\(\alpha+\beta\) は \(C\) の上限である。
	
	よって
	\[
	\sup C=\sup A+\sup B
	\]
	である。
\end{strategy}